% !TeX root = ../navier-stokes-manuscript.tex
% ============================================================
% 1.5 Class Membership and Participation
% ============================================================
\subsection{Class Membership and Participation}\label{sub:ClassMembershipAndParticipation}

Silver ends with a practical difficulty.
Every proposed loss of smoothness has to be tested against the same original fluid, so we need an exact way to say whether the proposed object still belongs to that fluid at all.

Fix the initial datum \(u_0\) on \(\mathbb R^3\) and one viscosity \(\nu>0\).
Let \((u,p)\) be the classical solution generated by those choices on its maximal interval \([0,T_*)\).
The object under discussion is always a proposed extension \(Q=(\widetilde u,\widetilde p)\) of this particular preterminal solution at \(T_*\).
The scope is narrow on purpose: the question concerns this evolution and what it can become, rather than every weak solution that might be constructed from unrelated data.

Write \(\operatorname{Member}(Q)\) when the proposed extension agrees with \((u,p)\) before \(T_*\) and continues to obey the requirements fixed at the beginning:
\[
  \partial_t\widetilde u
  +(\widetilde u\cdot\nabla)\widetilde u
  +\nabla\widetilde p
  =
  \nu\Delta\widetilde u,
  \qquad
  \nabla\cdot\widetilde u=0,
\]
with the same datum, viscosity, decay, energy condition, and pressure determined by the same velocity field.
Smoothness remains the conclusion to be proved.
It is not inserted into the meaning of \(\operatorname{Member}(Q)\).

Class exit is simply the complement
\[
  \operatorname{Exit}(Q)
  :=
  \neg\operatorname{Member}(Q).
\]
Gold and Silver can now be written against the same boundary.
Gold seeks
\[
  \operatorname{Member}(Q)
  \Longrightarrow
  \operatorname{Smooth}(Q),
\]
while Silver seeks the contrapositive
\[
  \neg\operatorname{Smooth}(Q)
  \Longrightarrow
  \operatorname{Exit}(Q).
\]
The two implications have the same logical content.
That equivalence fixes the direction of the Silver argument; the class-exit implication itself still has to be proved.
A large derivative, an unusual geometry, or a suggestive limiting picture gives evidence only after it identifies a requirement that the proposed extension can no longer satisfy.

\divider{What Participation Means}

To see what those requirements mean physically, follow one material parcel through the field.
If \(X(a,t)\) is the trajectory that begins at the material label \(a\), then
\[
  \frac{d}{dt}X(a,t)=u(t,X(a,t)).
\]
The parcel is carried into a new neighbourhood by the velocity itself.
Its velocity along that path changes according to
\[
  \frac{d}{dt}u(t,X(a,t))
  =
  D_tu(t,X(a,t))
  =
  -\nabla p(t,X(a,t))
  +\nu\Delta u(t,X(a,t)),
\]
where \(D_t=\partial_t+u\cdot\nabla\) is the material derivative.

This one line contains the physical relation we need.
The parcel moves into the velocity field around it, the nearby curvature of that field supplies the viscous part of its acceleration, and the pressure gradient comes from the incompressibility constraint imposed across the whole fluid.
As the parcel moves, its neighbours change.
As the surrounding velocity changes, both the local viscous response and the global pressure response change with it.
The parcel never acquires an independent equation of motion detached from that field.

This continuing dependence is what we mean by \emph{participation}.
A parcel participates when its motion remains part of the same coupled evolution: it is transported by \(u\), it changes velocity through the pressure and viscous terms generated by that \(u\), and it remains inside the divergence-free motion of the whole fluid.
A small region participates when the same relation holds throughout the region and agrees with the surrounding field along its boundary.

Participation allows perfectly ordinary zeroes.
A parcel can be instantaneously motionless when the pressure and viscous forces balance there.
A steady solution can have \(D_tu=0\) throughout space.
The zero solution has every derivative equal to zero and still satisfies the complete relation.
The decisive issue is whether the value at the parcel continues to agree with the field that determines it.

Couette flow gives one part of this relation a simple physical picture.
Place fluid between the plates \(x_2=0\) and \(x_2=h\), hold the lower plate fixed, and move the upper plate with constant tangential velocity \(\gamma h e_1\).
The steady interior profile is
\[
  u(x)=(\gamma x_2,0,0),
  \qquad
  p=\text{constant}.
\]
Its shear per unit density is
\[
  \tau_{12}=\nu\partial_{x_2}u_1=\nu\gamma.
\]
The nonzero shear expresses the transfer of tangential momentum from one layer to the next.

The caveat matters.
This linear profile also satisfies
\[
  \Delta u=0,
  \qquad
  (u\cdot\nabla)u=0.
\]
The viscous stresses on the two sides of an interior layer balance, so the bulk viscous force and the material acceleration are both zero there.
The moving plate continuously supplies the momentum that maintains the profile.
Couette flow isolates viscous shear between neighbouring layers.
Its moving boundaries keep it separate from an unforced whole-space evolution, and its zero Laplacian shows why a nonzero gradient alone does not force a nonzero acceleration.

In a general flow the local balance changes from one parcel to the next, transport continually changes which parcels are adjacent, and pressure keeps the whole motion incompressible at the same time.
An admitted smooth continuation participates in this complete relation.
A concrete failure of one of its terms can witness class exit, while loss of continuity or loss of finite control requires its own endpoint test.
